\section{MAKING THE BOUND OPERATIONAL: LEARNING ENOUGH ABOUT DETECTION TO USE IT}\label{sec:operational}

The bound in equation~\eqref{eq:mbar} does work only when the organization can
say something about $M$. Two things need to be clear about $M$ before any
route to it is described.

The first is what $M$ is not. $M = d_1/d_0$ is the ratio of two
probabilities: that an existing condition is detected after a prior finding,
and that it is detected after a clean history. It is not audit effort. Hours,
sample sizes, coverage percentages and procedure counts are not $M$, and
doubling any of them does not double $M$; they inform $M$ only through a
defensible connection to detection.

The second is that $M$ is target-specific. The multiplier is defined for a
particular detection target: a defined condition, finding class, procedure or
audit objective. Different compliance requirements, finding types, control
deficiencies and transaction tests can have different $d_0$, different $d_1$
and therefore different $M$. An organization should not assume that follow-up
has the same detection effect for every condition, requirement or procedure.
In practice the useful object is $M$ for a stated finding class, or a
defensible bound that applies to a specified family of conditions; the text
keeps the single symbol $M$ for readability, with that scope understood.

Six ways to learn or bound $M$ follow, each with different information
requirements and identifying assumptions. Table~\ref{tab:designs} summarizes
what each supplies and the assumption on which it turns.

\subsection*{Reference observation under a common measurement rule}

The cleanest design is reference observation under a common measurement rule:
a set of program-years in each history group, examined with the same
reference procedure, chosen so that the reference observations represent the
population whose condition rate is being estimated. The reference sample does
not have to ignore prior history. It may be stratified on prior-finding
status; what matters is that within each history group it represents that
group, through probability sampling or known weights. Tax administrations
maintain random examination programs beside operationally selected ones for
this reason \citep{slemrod2019}, and safety regulators have randomized
inspection allocation \citep{levine2012}.

Let $s$ be the sensitivity of the reference procedure for the target
condition, the probability it detects that condition when present. One extra
symbol is needed because the reference procedure is a different procedure
from the targeted audit, with its own detection probability, and $s$ is
unknown. If $s$ is the same in both history groups, the reference finding
rates are $q^{\mathrm{ref}}_0 = r_0\, s$ and $q^{\mathrm{ref}}_1 = r_1\, s$,
so
%
\begin{equation}\label{eq:ref}
  \frac{q^{\mathrm{ref}}_1}{q^{\mathrm{ref}}_0}
  \;=\;
  \frac{r_1\, s}{r_0\, s}
  \;=\;
  \Rlat .
\end{equation}
%
The reference procedure does not need perfect detection to recover the
underlying risk ratio. It needs the same sensitivity in both history groups,
because a common sensitivity cancels, and only then.
\begin{center}
\fbox{\parbox{0.9\textwidth}{\emph{A reference procedure does not need
perfect detection to identify the underlying risk ratio. It needs comparable
detection in the two history groups, and negligible false positives.}}}
\end{center}
With $\Rlat$ in hand, equation~\eqref{eq:product} gives the multiplier
directly:
%
\begin{equation}\label{eq:mfromref}
  M \;=\; \Robs / \Rlat .
\end{equation}
%
In words: the targeted audit gives the ratio of finding rates between the two
history groups. The reference work gives the ratio of underlying condition
rates, under a common measurement rule. Their ratio identifies the
multiplicative detection component $M$. If the targeted record shows the
fourfold ratio of Table~\ref{tab:worlds} and the reference sample shows
findings twice as often after a prior finding as after a clean year, then
$\Rlat = 2$ and $M = 4/2 = 2$: the observed fourfold ratio factorizes into a
twofold underlying-condition ratio and a twofold detection multiplier.
Section~\ref{sec:fac} could not do this because the Clearinghouse contains no
such sample.

Five assumptions carry the result (Proposition~\ref{prop:ref}). First, within
each prior-history group the reference sample represents the population whose
condition rate is being estimated, directly or through known sampling
weights. Second, the same reference measurement rule is applied in both
groups. Third, the reference procedure has common sensitivity across the two
groups for the target condition. Fourth, the reference observation does not
alter the condition before it is measured. Fifth, the reference procedure has
negligible false positives for the target condition, or its false-positive
process is separately modeled; a reference that records findings where no
condition exists does not cancel in the ratio. The third is the load-bearing
one,
and applying the same nominal procedure does not by itself guarantee it.
Sensitivity can differ between the groups if the conditions in one group are
more severe or of different types, if records are harder to obtain, if the
procedure is implemented differently, or if the case mix differs; then the
reference ratio is $\Rlat$ times the ratio of sensitivities and the
cancellation fails. Reference-sample size should be planned from the desired
precision and the expected condition or finding rates in the two history
groups.

\subsection*{High-sensitivity reference procedures}

A second anchor comes from procedures whose sensitivity for a narrowly defined
condition is known or credibly near-perfect. Where such a procedure is applied
to a whole subpopulation, its finding rate estimates that subpopulation's
condition rate, so $r$ is observed rather than bounded, and the targeted
procedure's finding rate in the same group then gives its detection
probability, $d = q/r$; anchors in both history groups yield $d_0$, $d_1$ and
$M$.

The anchor rests on the procedure's sensitivity, not on its coverage.
Exhaustive coverage by itself does not make detection perfect. It removes
sampling or coverage uncertainty, and it anchors the underlying rate only when
the procedure applied to that population has known or credibly near-perfect
sensitivity for the target condition. Where auditors speak of ``certainty
testing'', the certainty refers to inclusion in the tested population, not to
certainty that every existing condition will be recognized. Testing every
transaction for an exact duplicate payment may plausibly have very high
sensitivity for that narrowly defined condition; testing every file for a
judgment-dependent compliance failure may not, and the anchor is then only as
good as the sensitivity that can be defended.

\subsection*{Independent reference adjudication}

A representative or appropriately weighted sample of program-years within
each history group can be re-adjudicated by a reviewer who does not see the
original finding history and applies a fixed program of work. Blindness
removes the history's influence on the reviewer; it does not by itself create
ground truth, and two cases have to be kept apart.

In the first, a sufficiently accurate reference or adjudication procedure is
available for the target condition. Then the reference review approximates
latent status, $d_0$ and $d_1$ are estimated directly as the share of
reference-identified conditions the original audit had recorded in each
history group, and $M$ is their ratio; this is the audit analogue of
correcting a diagnostic study for differential verification
\citep{degroot2011}. In the second, the reference review is itself imperfect.
Then its own error process has to be addressed, through multiple independent
procedures, a model of detection that treats the reference as one more
imperfect measurement, or an external benchmark with known properties.
Ordinary blind re-performance of the same work does not automatically solve
the problem, because both reviews may miss the same condition for the same
reason, and the shared misses are invisible to both.

\subsection*{Two distinct detection procedures}

Two imperfect procedures can reveal how much each is missing. Suppose
procedure A vouches transaction-level evidence and procedure B obtains an
independent external confirmation or performs a genuinely different
control-based test, both applied to the same programs. Under conditional independence of the two procedures' failures and stable
sensitivity, and with no false positives, the share of the conditions B finds
that A also finds identifies A's detection probability; the share of the
conditions A finds that B also finds identifies B's; and the overlap then
says how many conditions both missed. This is the logic ecologists use to
estimate how often a species is present at sites where a survey did not see it
\citep{mackenzie2002}, and it needs no reference sample
(Proposition~\ref{prop:twoproc}).

Two requirements carry it. The procedures' failures must be unrelated once
the true condition is fixed, and each procedure's sensitivity must be stable
across the program-years being pooled. Using two operationally different
procedures is not enough on its own: if both tend to miss the same difficult
cases, for example where documentation does not exist, their agreement looks
reassuring while saying little about what both missed. Vouching paired with
an external confirmation can satisfy the first requirement; two reviewers
re-performing the same sample with the same method generally cannot.

\subsection*{Randomized or quasi-random changes in audit attention}

Some changes in how hard an audit looks arrive for reasons unrelated to the
underlying condition: randomized rotation, random reference selections,
externally imposed coverage rules, staffing changes plausibly unrelated to the
control environment, randomized review depth. These reveal how finding rates
change when observation changes. A change in audit attention is informative
about detection only if the change does not independently alter, or select
on, the underlying condition; that is the exclusion restriction, and a
staffing increase introduced because controls deteriorated fails it. And such
variation identifies relative detection effects between the regimes it
creates, or tightens bounds; absolute condition prevalence generally requires
an additional sensitivity anchor of the kind described above.

\subsection*{Procedure-based bounds}

An organization that can do none of the above may still be able to defend an
upper bound $\Mbar$ from what it knows about its own procedures: the fraction
of the population covered, sample sizes, the number and type of compliance
requirements tested, the use of targeted versus broad procedures, the strata
tested with a high-sensitivity procedure, historical quality-review results,
and any external validation study. These quantities are not $M$. They inform
$M$ only through an explicit model connecting audit procedures to detection
probability, for instance a statement of what share of existing conditions a
procedure of a given depth is expected to catch, and the bound is only as
defensible as that model.

The arithmetic once $\Mbar$ is in hand is simple. If the organization can
defend $M \le 1.5$ for the finding class in question and observes $\Robs =
3$, equation~\eqref{eq:mbar} gives $\Rlat \ge 3/1.5 = 2$: even after giving
follow-up the largest detection advantage the organization considers
defensible, the underlying condition must still be at least twice as common
in the prior-finding group, and with monotone detection the underlying ratio
lies between 2 and 3. What the record could only place between $q_1$ and
$1/q_0$ has become a range an audit committee can act on.

The remaining route is a structural model whose restrictions on how the
condition evolves and how detection responds are defended on audit grounds;
identification then has to be established for that specific model, and this
paper does not take that route.

\begin{table}[tbp]
\centering
\caption{\textbf{Six ways to learn or bound the follow-up detection
multiplier.} What each design supplies and the assumption on which it turns.
$M = d_1/d_0$ is defined for a stated condition, finding class or procedure;
the designs are not ranked, because they answer different questions with
different information requirements. Details and the remaining assumptions
for each are in Section~\ref{sec:operational}; the first two identifying
results are Propositions~\ref{prop:ref} and \ref{prop:twoproc} in
Appendix~\ref{app:id}.}
\label{tab:designs}
\small
\begin{tabular}{p{0.30\textwidth}p{0.30\textwidth}p{0.32\textwidth}}
\toprule
Design & What it supplies & Critical assumption \\
\midrule
Reference observation under a common measurement rule & $\Rlat$, hence $M = \Robs/\Rlat$ & Common sensitivity in the two history groups; no false positives \\
High-sensitivity reference procedure & Anchor for $r$, hence $d = q/r$ & Known or credibly near-perfect sensitivity for the target condition \\
Independent reference adjudication & $d_0$, $d_1$ directly & Reference status credible enough to stand in for the condition \\
Two distinct procedures & $r$, $d$ within a history group & Conditionally independent failures; stable sensitivity; no false positives \\
Randomized or quasi-random attention shift & Relative detection; tighter bounds & Exclusion restriction: the shift does not alter or select on the condition \\
Procedure-based bound & $\Mbar$, hence $\Rlat \ge \Robs/\Mbar$ & A defensible model linking procedures to detection probability \\
\bottomrule
\end{tabular}
\end{table}

